\corrigesChapitre{3}{Transformations chimiques et équations}
\par\noindent\begin{minipage}{\linewidth}
\begin{corrige}{1}{Fondre ou réagir}
\begin{questions}
\item Non : l’eau change d’état.
\item De nouvelles espèces, présentes dans la rouille, se forment en consommant notamment le fer.
\end{questions}
\end{corrige}
\end{minipage}
\par\noindent\begin{minipage}{\linewidth}
\begin{corrige}{2}{Avant et après}
\begin{questions}
\item Carbone et dioxygène.
\item Dioxyde de carbone.
\end{questions}
\end{corrige}
\end{minipage}
\par\noindent\begin{minipage}{\linewidth}
\begin{corrige}{3}{Un test fourni}
\begin{questions}
\item Le dioxyde de carbone.
\item Non ; plusieurs gaz et mécanismes peuvent produire des bulles.
\end{questions}
\end{corrige}
\end{minipage}
\par\noindent\begin{minipage}{\linewidth}
\begin{corrige}{4}{Compter}
\begin{questions}
\item Quatre H de chaque côté.
\item Deux O de chaque côté.
\end{questions}
\pcFigureEx{figures/cor-reaction-eau.pdf}{Les coefficients rétablissent la conservation des atomes ; les formules des espèces restent inchangées.}
\end{corrige}
\end{minipage}
\par\noindent\begin{minipage}{\linewidth}
\begin{corrige}{5}{Masse totale}
\begin{questions}
\item $44\unite g$.
\item Non : ce sont les atomes de chaque élément qui se conservent.
\end{questions}
\end{corrige}
\end{minipage}
\par\noindent\begin{minipage}{\linewidth}
\begin{corrige}{6}{Un gaz s’échappe}
\begin{questions}
\item $2{,}2\unite g$.
\item Non : elle est sortie, par exemple sous forme gazeuse ; il faut l’inclure dans le bilan complet.
\end{questions}
\pcFigureEx{figures/rev-cor-masse-gaz.pdf}{Le système pesé est ouvert : la différence correspond au gaz sorti avec les hypothèses données.}
\end{corrige}
\end{minipage}
\par\noindent\begin{minipage}{\linewidth}
\begin{corrige}{7}{Une masse inconnue}
\begin{questions}
\item $m=22-6=16\unite g$.
\item Un réactif restant devrait figurer dans la masse finale totale.
\end{questions}
\pcFigureEx{figures/rev-cor-masse-carbone.pdf}{Les deux réactifs consommés contribuent à la masse de dioxyde de carbone.}
\end{corrige}
\end{minipage}
\par\noindent\begin{minipage}{\linewidth}
\begin{corrige}{8}{Le bon inventaire}
\begin{questions}
\item Deux C, six H et sept O de chaque côté.
\item Non, ce sont des proportions d’entités dans le modèle de réaction.
\end{questions}
\par\smallskip\noindent\textbf{Pour comprendre.}\enspace À gauche : C = 2, H = 6 et O = $1+3\times2=7$. À droite : C = $2\times1=2$, H = $3\times2=6$ et O = $2\times2+3=7$. Les nombres d’atomes sont égaux pour chaque élément.\par
\pcFigureEx{figures/rev-cor-ethanol.pdf}{Inventaire atomique de chaque côté ; l’oxygène de l’éthanol fait aussi partie du bilan.}
\end{corrige}
\end{minipage}
\par\noindent\begin{minipage}{\linewidth}
\begin{corrige}{9}{Une équation impossible}
\begin{questions}
\item Elle compte deux O à gauche et un à droite.
\item La seconde conserve quatre H et deux O de chaque côté.
\end{questions}
\pcFigureEx{figures/cor-reaction-eau.pdf}{Les coefficients rétablissent la conservation des atomes ; les formules des espèces restent inchangées.}
\end{corrige}
\end{minipage}
\par\noindent\begin{minipage}{\linewidth}
\begin{corrige}{10}{Des réactifs en excès}
\begin{questions}
\item $10+25-5=30\unite g$.
\item Les 35 g finaux comprennent aussi les 5 g non transformés.
\end{questions}
\par\smallskip\noindent\textbf{Pour comprendre.}\enspace La masse initiale totale vaut $10+25=35\unite g$. S’il reste 5 g de réactif, les produits formés représentent $35-5=30\unite g$ ; la masse finale totale reste 35 g.\par
\pcFigureEx{figures/rev-cor-reactif-restant.pdf}{Ne pas confondre masse des produits formés et masse finale totale avec le réactif restant.}
\end{corrige}
\end{minipage}
\par\noindent\begin{minipage}{\linewidth}
\begin{corrige}{11}{Comparer des contrôles}
\begin{questions}
\item Le changement peut venir du réactif, du milieu ou d’une contamination.
\item Non, il faut contrôler la cause et reprendre un protocole interprétable.
\end{questions}
\pcFigureEx{figures/rev-cor-temoin.pdf}{Le changement du témoin empêche d’attribuer avec certitude le résultat au gaz étudié.}
\end{corrige}
\end{minipage}
\par\noindent\begin{minipage}{\linewidth}
\begin{corrige}{12}{Mélange et réaction}
\begin{questions}
\item Non : les molécules de sucre restent présentes, dispersées dans le solvant.
\item Des substances mises en contact peuvent réagir ; il faut examiner les espèces initiales et finales.
\end{questions}
\end{corrige}
\end{minipage}
\par\noindent\begin{minipage}{\linewidth}
\begin{corrige}{13}{Un solide gagne de la masse}
\begin{questions}
\item Le dioxygène incorporé au solide, ici 0,8 g.
\item $2{,}1+0{,}8=2{,}9\unite g$.
\end{questions}
\pcFigureEx{figures/rev-cor-masse-fer.pdf}{L’augmentation de masse vient de matière du milieu incorporée dans le produit.}
\end{corrige}
\end{minipage}
\par\noindent\begin{minipage}{\linewidth}
\begin{corrige}{14}{Un protocole de gaz}
\begin{questions}
\item Identifier le gaz avec une banque de tests, comparer les espèces initiales et finales, et prévoir un témoin.
\item La quantité peut être trop faible ou le réactif inadéquat ; une absence de résultat visible a une portée limitée.
\end{questions}
\end{corrige}
\end{minipage}
\par\noindent\begin{minipage}{\linewidth}
\begin{corrige}{15}{Prise d’initiative}
\begin{questions}
\item Non : le bilan de masse est cohérent mais ne donne pas l’identité.
\item Un test caractéristique, une formule dans une équation fournie ou d’autres propriétés mesurées.
\end{questions}
\end{corrige}
\end{minipage}
\par\noindent\begin{minipage}{\linewidth}
\phantomsection\label{cor-defi-3}\pcSolutionDefi{3}{Le bilan complet}\begin{questions}
\item $m=8+40-8-22=18\unite g$.
\item Avant : 48 g ; après : $8+22+18=48\unite g$.
\item L’enceinte est un contenant supposé inerte. Sa masse intervient dans la pesée totale, mais elle est inchangée et ne représente pas une espèce consommée.
\end{questions}
\par\smallskip\noindent\textbf{Pour comprendre.}\enspace La masse initiale est $8+40=48\unite g$. La masse finale est $22+18+8=48\unite g$. Les 8 g de dioxygène restant doivent figurer dans le bilan ; ils ne font pas partie des produits formés.\par
\pcFigureEx{figures/rev-cor-masse-methane.pdf}{Bilan complet du modèle : réactifs consommés, produits et dioxygène restant.}
\end{minipage}
\par\noindent\begin{minipage}{\linewidth}
\phantomsection\label{cor-activite-3}\pcSousTitre{Pistes pour l’investigation}
Fusion et dissolution sans réaction ne forment pas de nouvelles espèces ; un dégazage peut seulement libérer une espèce déjà présente. Une réaction nécessite une interprétation des espèces consommées et formées. Le test et son témoin renforcent l’argument, sans que les seules bulles suffisent. Identifier le système et les échanges avant tout bilan de masse.\par\end{minipage}
